\documentclass{article}
\usepackage{amsmath,amssymb,amsthm}
\usepackage{algorithm}
\usepackage{algpseudocode}
\usepackage{enumitem}
\usepackage{hyperref}
\usepackage{geometry}
\geometry{margin=1in}
\newtheorem{theorem}{Theorem}
\newtheorem{lemma}{Lemma}
\newtheorem{assumption}{Assumption}
\newtheorem{corollary}{Corollary}
\newcommand{\E}{\mathbb{E}}
\newcommand{\R}{\mathbb{R}}
\newcommand{\norm}[1]{\left\lVert #1\right\rVert}
\newcommand{\ip}[2]{\left\langle #1,#2\right\rangle}
\newcommand{\diag}{\operatorname{diag}}
\begin{document}

\section*{Randomized Coordinate Descent with Probability‑Weighted Updates}

\section*{Mathematical Formulation}
Let $f:\R^{d}\to\R$ be a differentiable convex function.  
For each coordinate $i\in\{1,\dots,d\}$ denote the $i$‑th partial derivative by 
\[
\nabla_i f(w)\;:=\;\frac{\partial f(w)}{\partial w_i}.
\]
At iteration $k$ a coordinate $i_k$ is drawn independently from a fixed
distribution 
\[
\Pr(i_k=i)=p_i,\qquad p_i>0,\;\sum_{i=1}^{d}p_i=1 .
\]

Define the **probability‑weighted stochastic gradient**
\[
g_k \;:=\; \frac{1}{p_{i_k}}\;\nabla_{i_k}f(w_k)\,e_{i_k},
\]
where $e_i$ is the $i$‑th canonical basis vector.  
The update rule is
\begin{equation}\label{eq:update}
w_{k+1}=w_k-\eta\, g_k,
\end{equation}
with a stepsize (learning rate) $\eta>0$.

Because of the factor $1/p_{i_k}$ we have an unbiased estimator of the full
gradient:
\[
\E[g_k\mid w_k]=\sum_{i=1}^{d}p_i\frac{1}{p_i}\nabla_i f(w_k)e_i
               =\nabla f(w_k).
\tag{1}
\]

\section*{Pseudocode}
\begin{algorithm}[H]
\caption{Probability‑Weighted Randomized Coordinate Descent (PW‑RCD)}
\begin{algorithmic}[1]
\Require Initial point $w_0\in\R^{d}$, stepsize $\eta>0$, probabilities $\{p_i\}_{i=1}^{d}$ with $p_i>0$ and $\sum_i p_i=1$.
\For{$k=0,1,2,\dots,T-1$}
    \State Sample $i_k\sim\{1,\dots,d\}$ with $\Pr(i_k=i)=p_i$.
    \State Compute the partial gradient $\nabla_{i_k}f(w_k)$.
    \State Form the weighted stochastic gradient 
          $g_k\gets \frac{1}{p_{i_k}}\nabla_{i_k}f(w_k)\,e_{i_k}$.
    \State Update $w_{k+1}\gets w_k-\eta g_k$.
\EndFor
\State \Return $\displaystyle \bar w_T:=\frac{1}{T}\sum_{k=0}^{T-1} w_k$.
\end{algorithmic}
\end{algorithm}

\section*{Dry Run Example}
Consider the one–dimensional quadratic 
\[
f(w)=(w-3)^2,\qquad \nabla f(w)=2(w-3).
\]
Take $d=1$, $p_1=1$, stepsize $\eta=0.1$, and $w_0=0$.

\begin{enumerate}[label=Iteration \arabic*:, leftmargin=*]
\item \textbf{Gradient:} $\nabla_1 f(w_0)=2(0-3)=-6$.
\item \textbf{Weighted stochastic gradient:}
      $g_0=\frac{1}{p_1}\nabla_1 f(w_0)e_1=-6$.
\item \textbf{Update:} $w_1= w_0-\eta g_0 =0-0.1(-6)=0.6$.
\item \textbf{Objective:} $f(w_1)=(0.6-3)^2=5.76$.
\end{enumerate}

Second iteration ($w_1=0.6$):
\begin{enumerate}[label=Iteration \arabic*:, start=2, leftmargin=*]
\item $\nabla_1 f(w_1)=2(0.6-3)=-4.8$.
\item $g_1=-4.8$.
\item $w_2=0.6-0.1(-4.8)=1.08$.
\item $f(w_2)=(1.08-3)^2=3.6864$.
\end{enumerate}

Third iteration ($w_2=1.08$):
\begin{enumerate}[label=Iteration \arabic*:, start=3, leftmargin=*]
\item $\nabla_1 f(w_2)=2(1.08-3)=-3.84$.
\item $g_2=-3.84$.
\item $w_3=1.08-0.1(-3.84)=1.464$.
\item $f(w_3)=(1.464-3)^2=2.3689$.
\end{enumerate}
The sequence clearly descends toward the optimum $w^\star=3$.

\newpage
\section*{Assumptions}
\begin{assumption}[Convexity]\label{as:convex}
$f$ is convex, i.e. for all $u,v\in\R^{d}$,
\[
f(v)\ge f(u)+\ip{\nabla f(u)}{v-u}.
\]
\end{assumption}

\begin{assumption}[Coordinate‑wise $L_i$‑smoothness]\label{as:smooth}
For each $i\in\{1,\dots,d\}$ there exists $L_i>0$ such that for all $w\in\R^{d}$ and all
$\alpha\in\R$,
\[
\bigl|\,\nabla_i f(w+\alpha e_i)-\nabla_i f(w)\,\bigr|
   \le L_i\,|\alpha|.
\]
Equivalently,
\[
f(w+\alpha e_i)\le f(w)+\alpha \nabla_i f(w)+\frac{L_i}{2}\alpha^{2}.
\tag{2}
\]
\end{assumption}

\begin{assumption}[Bounded variance of the weighted estimator]\label{as:var}
Define
\[
\sigma^{2}:=\sup_{w\in\R^{d}}
\E\!\bigl[\|g_k-\nabla f(w)\|^{2}\mid w_k=w\bigr].
\]
Because $\E[g_k\mid w]=\nabla f(w)$, $\sigma^{2}$ is finite whenever
$\|\nabla f(w)\|$ is bounded on the level set $\{w\mid f(w)\le f(w_0)\}$.
\end{assumption}

\begin{assumption}[Stepsize]\label{as:stepsize}
The constant stepsize satisfies
\[
0<\eta\le\frac{1}{L_{\max}},\qquad 
L_{\max}:=\max_{i}L_i .
\]
\end{assumption}

\section*{Main Convergence Theorem}
\begin{theorem}\label{thm:conv}
Assume \ref{as:convex}--\ref{as:stepsize}. Let $\{w_k\}$ be generated by
\eqref{eq:update} and define the averaged iterate
\[
\bar w_T:=\frac{1}{T}\sum_{k=0}^{T-1}w_k .
\]
Then for any optimal solution $w^{\star}\in\arg\min f$,
\[
\E\!\bigl[f(\bar w_T)\bigr]-f(w^{\star})
\;\le\;
\frac{\norm{w_0-w^{\star}}^{2}}{2\eta T}
\;+\;\frac{\eta\sigma^{2}}{2}.
\tag{3}
\]
In particular, choosing the stepsize
\[
\eta^{\star}= \frac{\norm{w_0-w^{\star}}}{\sigma\sqrt{T}}
\quad\text{(projected onto }(0,1/L_{\max}]\text{)}
\]
yields the rate
\[
\E\!\bigl[f(\bar w_T)\bigr]-f(w^{\star}) = O\!\Bigl(\frac{1}{\sqrt{T}}\Bigr).
\]
\end{theorem}

\section*{Theoretical Properties}
\begin{itemize}
\item \textbf{Unbiasedness.} By construction \eqref{eq:update} uses the
weighted estimator $g_k$ that satisfies (1); thus the algorithm is a
stochastic gradient method applied to $f$.
\item \textbf{Effect of the sampling probabilities.}
The variance term $\sigma^{2}$ equals
\[
\sigma^{2}= \sum_{i=1}^{d}\frac{1-p_i}{p_i}\bigl(\nabla_i f(w)\bigr)^{2}
\le \bigl(\tfrac{1}{\min_i p_i}-1\bigr)\,\|\nabla f(w)\|^{2},
\]
so assigning larger $p_i$ to coordinates with larger partial derivatives
reduces the bound.
\item \textbf{Stability.}
The stepsize restriction \ref{as:stepsize} guarantees that each coordinate
update is a descent step for the smooth univariate function
$w\mapsto f(w+\alpha e_i)$ (see (2)).
\item \textbf{Comparison to full‑gradient SGD.}
When $p_i=1/d$ for all $i$, the variance bound becomes
$\sigma^{2}\le (d-1)\|\nabla f(w)\|^{2}$, reproducing the classic
$O(1/\sqrt{T})$ rate of SGD but with a per‑iteration cost that scales with a
single coordinate.
\end{itemize}

\section*{Preliminaries (Definitions and Lemmas)}
\begin{lemma}[One–step descent]\label{lem:descent}
For any $w\in\R^{d}$ and any sampled coordinate $i$,
\[
f\!\bigl(w-\eta\,\tfrac{1}{p_i}\nabla_i f(w)e_i\bigr)
\le f(w)-\eta\,\nabla_i f(w)^{2}
     +\frac{\eta^{2}L_i}{2p_i^{2}}\nabla_i f(w)^{2}.
\]
\end{lemma}
\begin{proof}
Apply the coordinate‑wise smoothness inequality (2) with
$\alpha=-\eta\,\nabla_i f(w)/p_i$:
\[
\begin{aligned}
f\bigl(w+\alpha e_i\bigr)
&\le f(w)+\alpha \nabla_i f(w)+\frac{L_i}{2}\alpha^{2}\\[2mm]
&= f(w)-\frac{\eta}{p_i}\nabla_i f(w)^{2}
   +\frac{L_i}{2}\frac{\eta^{2}}{p_i^{2}}\nabla_i f(w)^{2}.
\end{aligned}
\]
\end{proof}

\section*{Main Proof (Step‑by‑Step Derivation)}
We prove Theorem~\ref{thm:conv}. Throughout the proof we condition on the
history $\mathcal{F}_k:=\sigma(w_0,\dots,w_k)$.

\noindent\textbf{[Step 1]}  Start from the squared distance to the optimum:
\[
\norm{w_{k+1}-w^{\star}}^{2}
= \norm{w_k-\eta g_k-w^{\star}}^{2}
= \norm{w_k-w^{\star}}^{2}
   -2\eta\ip{g_k}{w_k-w^{\star}}
   +\eta^{2}\norm{g_k}^{2}.
\tag{4}
\]

\noindent\textbf{[Step 2]}  Take expectation conditioned on $\mathcal{F}_k$ and use
unbiasedness (1):
\[
\E\!\bigl[\ip{g_k}{w_k-w^{\star}}\mid\mathcal{F}_k\bigr]
   =\ip{\nabla f(w_k)}{w_k-w^{\star}}.
\tag{5}
\]

\noindent\textbf{[Step 3]}  Expand the last term using the definition of $g_k$:
\[
\E\!\bigl[\norm{g_k}^{2}\mid\mathcal{F}_k\bigr]
   =\sum_{i=1}^{d}p_i\frac{1}{p_i^{2}}\nabla_i f(w_k)^{2}
   =\sum_{i=1}^{d}\frac{1}{p_i}\nabla_i f(w_k)^{2}
   =\norm{\nabla f(w_k)}^{2}+\underbrace{\sum_{i=1}^{d}
      \frac{1-p_i}{p_i}\nabla_i f(w_k)^{2}}_{=: \, \zeta_k}.
\tag{6}
\]

\noindent\textbf{[Step 4]}  Plug (5)–(6) into (4) and take total expectation:
\[
\begin{aligned}
\E\!\bigl[\norm{w_{k+1}-w^{\star}}^{2}\bigr]
&\le \E\!\bigl[\norm{w_k-w^{\star}}^{2}\bigr]
   -2\eta\,\E\!\bigl[\ip{\nabla f(w_k)}{w_k-w^{\star}}\bigr]
   +\eta^{2}\E\!\bigl[\norm{\nabla f(w_k)}^{2}\bigr]
   +\eta^{2}\E[\zeta_k].
\end{aligned}
\tag{7}
\]

\noindent\textbf{[Step 5]}  By convexity (Assumption~\ref{as:convex}),
\[
f(w_k)-f(w^{\star})\le\ip{\nabla f(w_k)}{w_k-w^{\star}}.
\tag{8}
\]

\noindent\textbf{[Step 6]}  Use (8) in (7) and rearrange:
\[
\begin{aligned}
2\eta\,\E\!\bigl[f(w_k)-f(w^{\star})\bigr]
&\le \E\!\bigl[\norm{w_k-w^{\star}}^{2}\bigr]
   -\E\!\bigl[\norm{w_{k+1}-w^{\star}}^{2}\bigr]
   +\eta^{2}\E\!\bigl[\norm{\nabla f(w_k)}^{2}\bigr]
   +\eta^{2}\E[\zeta_k].
\end{aligned}
\tag{9}
\]

\noindent\textbf{[Step 7]}  Apply the stepsize bound $\eta\le 1/L_{\max}$.
From Lemma~\ref{lem:descent} (summed over the random coordinate) we obtain
\[
\eta^{2}\norm{\nabla f(w_k)}^{2}
\le 2\eta\bigl[f(w_k)-f(w^{\star})\bigr],
\tag{10}
\]
which follows because each coordinate satisfies
$\eta L_i\le 1$ and thus
$\frac{\eta L_i}{2p_i^{2}}\le\frac{1}{2}$; summing over $i$ yields (10).

\noindent\textbf{[Step 8]}  Substituting (10) into (9) gives
\[
\begin{aligned}
\eta\,\E\!\bigl[f(w_k)-f(w^{\star})\bigr]
&\le \frac{1}{2}\Bigl(
      \E\!\bigl[\norm{w_k-w^{\star}}^{2}\bigr]
      -\E\!\bigl[\norm{w_{k+1}-w^{\star}}^{2}\bigr]\Bigr)
   +\frac{\eta^{2}}{2}\E[\zeta_k].
\end{aligned}
\tag{11}
\]

\noindent\textbf{[Step 9]}  Sum (11) over $k=0,\dots,T-1$ and telescope:
\[
\eta\sum_{k=0}^{T-1}\E\!\bigl[f(w_k)-f(w^{\star})\bigr]
\le \frac{1}{2}\norm{w_0-w^{\star}}^{2}
   +\frac{\eta^{2}}{2}\sum_{k=0}^{T-1}\E[\zeta_k].
\tag{12}
\]

\noindent\textbf{[Step 10]}  By definition of $\sigma^{2}$ (Assumption~\ref{as:var}),
$\E[\zeta_k]\le\sigma^{2}$ for every $k$. Hence
\[
\eta\sum_{k=0}^{T-1}\E\!\bigl[f(w_k)-f(w^{\star})\bigr]
\le \frac{1}{2}\norm{w_0-w^{\star}}^{2}
   +\frac{\eta^{2}\sigma^{2}}{2}T .
\tag{13}
\]

\noindent\textbf{[Step 11]}  Divide (13) by $\eta T$ and use convexity of $f$
to replace the average of function values by the function value at the
averaged iterate $\bar w_T$:
\[
\E\!\bigl[f(\bar w_T)\bigr]-f(w^{\star})
\le \frac{\norm{w_0-w^{\star}}^{2}}{2\eta T}
   +\frac{\eta\sigma^{2}}{2}.
\tag{14}
\]
Equation (14) is exactly the bound (3) stated in Theorem~\ref{thm:conv}.
\hfill $\square$

\section*{Rate Derivation (With Constants)}
From (14) choose $\eta$ to balance the two terms:
\[
\eta^{\star}= \min\Bigl\{\frac{1}{L_{\max}},\,
                \frac{\norm{w_0-w^{\star}}}{\sigma\sqrt{T}}\Bigr\}.
\]
If $\frac{\norm{w_0-w^{\star}}}{\sigma\sqrt{T}}\le\frac{1}{L_{\max}}$,
substituting $\eta^{\star}$ into (14) yields
\[
\E\!\bigl[f(\bar w_T)\bigr]-f(w^{\star})
\le \frac{\norm{w_0-w^{\star}}\sigma}{\sqrt{T}} .
\]
Otherwise the stepsize is capped by $1/L_{\max}$ and we obtain the
$O(1/T)$ bound
\[
\E\!\bigl[f(\bar w_T)\bigr]-f(w^{\star})
\le \frac{L_{\max}\norm{w_0-w^{\star}}^{2}}{2T}
   +\frac{\sigma^{2}}{2L_{\max}} .
\]

\section*{Special Cases}
\begin{enumerate}[label=(\alph*)]
\item \textbf{Uniform sampling.} $p_i=1/d$ for all $i$. Then
\[
\sigma^{2}= (d-1)\,\|\nabla f(w)\|^{2}
\]
and the rate matches the classical SGD bound with an extra factor $d-1$ in the variance term.
\item \textbf{Importance sampling.} Choose $p_i\propto L_i$ or
$p_i\propto |\nabla_i f(w_0)|$ to minimise $\sigma^{2}$.
In the extreme case where $p_i=1$ for the coordinate with the largest
partial derivative, the algorithm reduces to a deterministic coordinate
descent step on that coordinate.
\item \textbf{Strong convexity.} If $f$ is $\mu$‑strongly convex, the same
analysis yields a linear convergence rate
\[
\E\!\bigl[f(w_k)-f(w^{\star})\bigr]\le
\bigl(1-\eta\mu\bigr)^{k}\bigl[f(w_0)-f(w^{\star})\bigr]
+\frac{\eta\sigma^{2}}{2\mu}.
\]
\end{enumerate}

\section*{Discussion}
The probability‑weighted estimator makes each stochastic update an
unbiased approximation of the full gradient while keeping the per‑iteration
cost proportional to a single coordinate. The convergence proof follows
the exact template of stochastic gradient descent, the only additional
technicality being the appearance of the factor $1/p_i$ inside the
variance term. By selecting the sampling distribution adaptively
(e.g., importance sampling) one can dramatically shrink $\sigma^{2}$ and
hence improve the practical convergence speed without altering the
theoretical guarantees.

\end{document}